% ============================================================
% Paper 4 — A Dual Operator for Prime–Zero Coupling
%            and a Conditional Proof of Energy Asymmetry
%
% Author:  Ulrich Tehrani
% Version: v31 (August 2026)
% DOI:     10.5281/zenodo.19364703
% License: CC BY 4.0
% ============================================================
\documentclass[11pt]{article}
\usepackage{cmap}            % ToUnicode CMaps for §, en/em-dashes etc.
\usepackage[T1]{fontenc}
\emergencystretch=1em        % avoid overfull hboxes from inline math
\usepackage[utf8]{inputenc}
\usepackage[a4paper,margin=1in]{geometry}
\usepackage{amsmath,amssymb,amsthm,mathtools}
\usepackage{xcolor}
\usepackage{hyperref}
\pdfstringdefDisableCommands{%
  \def\Tt{T-tilde}%
  \def\etaren{eta-ren}%
  \def\etaorig{eta-orig}%
}
\usepackage{enumitem}
\usepackage{setspace}
\usepackage{booktabs}
\usepackage{array}
\usepackage{graphicx}

\hypersetup{colorlinks=true,linkcolor=blue!50!black,urlcolor=blue!50!black,citecolor=blue!50!black}
\setlength{\parskip}{0.5em}
\setlength{\parindent}{0pt}
\onehalfspacing

% -- Theorem environments
\newtheorem{theorem}{Theorem}[section]
\newtheorem{proposition}[theorem]{Proposition}
\newtheorem{lemma}[theorem]{Lemma}
\newtheorem{corollary}[theorem]{Corollary}
\newtheorem{definition}[theorem]{Definition}
\newtheorem{assumption}[theorem]{Assumption}
\newtheorem{observation}[theorem]{Observation}
\newtheorem*{remark}{Remark}
\newtheorem{openproblem}[theorem]{Open Problem}

% Custom label-free environments
\newenvironment{model}[1]{\par\smallskip\noindent\textit{[Model: #1]}\quad}{\par\smallskip}
\newenvironment{heuristic}[1]{\par\smallskip\noindent\textit{[Heuristic: #1]}\quad}{\par\smallskip}

% -- Macros (Paper 3 base)
\newcommand{\Hstr}{H_{\mathrm{str}}}
\newcommand{\Hnull}{H_{\mathrm{null}}}
\newcommand{\Estr}{E_{\mathrm{str}}}
\newcommand{\Espec}{E_{\mathrm{spec}}}
\newcommand{\etaorig}{\eta_{\mathrm{orig}}}
\newcommand{\etaren}{\eta_{\mathrm{ren}}}
\newcommand{\Tren}{T_{\mathrm{ren}}}
\newcommand{\Gun}{G^{\mathrm{un}}}
\newcommand{\Hlocal}{H_{\mathrm{local}}}
% -- Paper 4 specific macros
\newcommand{\Tt}{\widetilde{T}}
\newcommand{\fp}{f_p}
\newcommand{\Ceta}{C_{\eta}}
\newcommand{\CT}{C_{T}}
\newcommand{\kdom}{k_{\mathrm{dom}}}
\newcommand{\Wone}{W_1}
\newcommand{\Dburstvar}{\Delta_{\mathrm{Burst}}}
\newcommand{\Dcross}{\Delta_{\mathrm{Cross}}}
\newcommand{\Dstream}{\Delta_{\mathrm{Stream}}}
\newcommand{\Da}{D_a}

% -- Title
\title{\vspace{-1.5em}\bfseries A Dual Operator for Prime--Zero Coupling\\[0.4em]
\large and a Conditional Proof of Energy Asymmetry}
\author{Ulrich Tehrani}
\date{Research Note \textperiodcentered\ August 2026 \textperiodcentered\ v31}

\begin{document}
\maketitle

% -------------------------------------------------------
\begin{abstract}
We introduce the operator $\Tt := \Phi \circ \Phi^*$,
the dual of the loop operator $T = \Phi^* \circ \Phi$ constructed in Paper~3
of this series. Here $\Phi\colon\Hstr \to \Hnull$ maps a
finite-dimensional prime space into a finite-dimensional zero space via
Gaussian-damped prime resonance vectors.

\textbf{What is proved.}
We establish (\S\,2) the algebraic spectral identity
$\sigma(T)\setminus\{0\} = \sigma(\Tt)\setminus\{0\}$
by classical operator theory.
We prove (\S\,3.4) that $\Wone := \CT\cdot\Tt^+$ is self-adjoint.
$\Wone$ does not assume RH or the HP postulate; however, its normalization
constant $\CT$ is calibrated against the zero ordinates via OLS and should
not be presented as fully input-free.
$\Tt$ is explicit and deterministic; the $\gamma_k$ are external numerical
inputs, not conclusions of the model (\S\,3.4).
We prove (\S\,4) the Abel Summation Principle (Lemma~\ref{lem:M3}):
for monotone decreasing weights $f_i$ and a sequence $g_i$ with bounded
partial sums, $|\sum_i f_i g_i| \leq M f_1$.
We prove (\S\,4) the Prime Exponential Endpoint Estimate
(Proposition~\ref{thm:EXPLICIT}): for each fixed zero ordinate $\gamma_k$,
\[
M_k(\kappa) := \max_{x \leq \kappa}\Bigl|\sum_{p \leq x}\sin(\gamma_k\log p)\Bigr|
\le \left(\frac{1}{\sqrt{1+\gamma_k^2}}+o_{\gamma_k}(1)\right)\pi(\kappa),
\]
with no uniformity in $k$ asserted,
using the classical de~la~Vall\'ee-Poussin form of the Prime Number Theorem,
and no RH, GUE, or HP input.

\textbf{What is numerical.}
The algebraic spectral identity is verified numerically:
$\max_j|\lambda_j(T)-\mu_j(\Tt)|<10^{-15}$ [Numerical, \S\,2.3].
At the benchmark grid $(\kappa,\varepsilon,N)=(53,0.05,100)$, the
leading tested eigenvectors ($j=1,\ldots,14$) of $\Tt$ localize at
zero ordinates with observed degrees $0.45$--$0.99$ [Numerical, \S\,3.2].
The correlation $r_1 = \mathrm{corr}(\mu_j,1/\gamma_{k(j)}) = 0.950$
($\kappa=53$) [Numerical, \S\,3.1].
The energy asymmetry $\etaren > 0$ holds on the tested grid
$\kappa\in\{23,53,101,199,503,1009\}$
[Numerical, \S\,5.1].
The operator $\Tt$ is not supported as a Hilbert--P\'olya operator
by the tested spectral diagnostics: $r_2 = \mathrm{corr}(\omega_j,\gamma_{k(j)})$
falls from $0.50$ to $0.16$ on the tested grid
$\kappa\in\{23,\ldots,503\}$ [Numerical, \S\,2.4].

\textbf{What is conditional.}
Under Assumption $(E_{\mathrm{rem}})$ (remainder control:
$|\Dcross+\Dstream|\le\rho\cdot\Dburstvar$ for some $\rho<1$,
numerically $\rho_{\max}=0.583<1$ on the full tested grid
$\kappa\in\{23,53,101,199,503,1009\}$) we prove
(\S\,5.3) that $\etaren(\kappa)\ge(1-\rho)\Dburstvar/E_{\mathrm{str}}(\kappa)>0$
[conditional, by triangle inequality].
We establish numerically that
$\Dburstvar = 3.105\ldots > 0$ is $\kappa$-invariant by definition for
$\kappa\ge 7$ at fixed $(\varepsilon,N)=(0.05,100)$.
On the tested grid $\kappa\in\{23,53,101,199,503,1009\}$,
one observes $\Delta(\kappa)\ge\Dburstvar$; analytically,
this stronger lower-bound property is OPEN.
Under Assumption $(E_{\mathrm{rem}})$, the conditional theorem gives
$\Delta(\kappa)\ge(1-\rho)\Dburstvar>0$,
and on the tested grid $\kappa\in\{23,53,101,199,503,1009\}$,
the values of $\etaren(\kappa)$ are positive [numerical extrapolation; no convergence theorem].

\textbf{What remains open.}
The analytic proof of $\etaren > 0$ without the above assumption
(\S\,6); whether a circularity-free spectral function
$f(\Tt)$ exists with spectrum converging to $\{\gamma_k\}$
(\S\,6); and the arithmetic origin of
$\Ceta \approx 0.39$ (\S\,6).

\textbf{Structure.}
The algebraic statements in \S\S\,2--3 are non-circular
and proof-complete; the spectral diagnostics in \S\,3 are numerical.
\S\,4 contains proved estimates; \S\,5 contains the
numerical and conditional energy-asymmetry analysis.
Section~6 states the open problems.
No use of the Riemann Hypothesis, the GUE conjecture, the Montgomery
pair correlation conjecture, or the Hilbert--P\'olya postulate is made
anywhere in \S\S\,2--5.
\end{abstract}

\vspace{2em}

% --------------------------------------------------------
% Notation and Conventions
% --------------------------------------------------------
\section*{Notation and Conventions}
\label{sec:notation}

\emph{Critical line.}
Throughout this series, the \emph{critical line} means
$\Re(s)=\tfrac12$, parameterised by $s=\tfrac12+i\gamma$
with $\gamma\in\mathbb{R}$.

\medskip
\emph{Global parameters.}
The benchmark parameter set is $\kappa=53$ ($\pi(\kappa)=16$ active primes),
$\varepsilon=0.05$, $N=100$ Riemann zero ordinates, $\sigma_0=\tfrac12$.
Numerical sections explicitly state when $\kappa$ is varied;
$\varepsilon=0.05$ and $N=100$ are fixed unless stated otherwise.

\medskip
\emph{Main objects.}
The encoding map $\Phi\colon\Hstr\to\Hnull$
assigns to each prime $p\leq\kappa$ a resonance vector
$a_p=(e^{-\varepsilon^2\gamma_k^2/2}\sin(\gamma_k\log p))_{k=1}^N$.
The loop operator $T=\Phi^*\!\circ\Phi$ acts on the prime
space $\Hstr$; its dual, the Tehrani operator
$\Tt=\Phi\circ\Phi^*$, acts on the zero space $\Hnull$
and is the central object of this paper.
The canonical weight $c_p^{\mathrm{ren}}:=\sqrt{f_p}$ is the
Weil-renormalised weight from Paper~2, \S\,4.
The energy asymmetry is
$\etaren := \eta[c^{\mathrm{ren}}] = 1-c^T\Gun c/\Estr$,
where $\eta[c]$ denotes the asymmetry functional evaluated at weight $c$.
Since all Paper-4 energy computations are at the fixed reference line
$\sigma_0=\tfrac12$, we write $\etaren(\kappa)$ for $\eta[c^{\mathrm{ren}}]$
at fixed $(\sigma_0,\varepsilon,N)=(\tfrac12,0.05,100)$.
This differs from Paper~3's $\etaorig := \eta[c^{\eta}]$, which uses
$c_p^{\eta}=\sqrt{V_p(\tfrac12)}$. At $\kappa=53$ and fixed
$(\sigma_0,\varepsilon,N)=(\tfrac12,0.05,100)$, the two reference values are
\[
\etaren=0.6693,\qquad \etaorig=0.6908.
\]
Two arithmetic constants govern spectral scaling:
$\Ceta\approx 0.39$ (renormalised eigenvalue ratio, Definition~\ref{def:constants}) and
$\CT$ (OLS slope of $\mu_j$ vs.\ $1/\gamma_{k(j)}$).

\medskip
\emph{Analytical foundations.}
The proofs use the Prime Number Theorem
(Hadamard 1896, de~la~Vall\'ee-Poussin 1896),
the Abel partial summation formula, standard properties of
Moore--Penrose pseudoinverses, and standard
finite-dimensional spectral theory, including the
singular-value identity for $\Phi^*\Phi$ and $\Phi\Phi^*$.

\medskip
\emph{Numerical computations.}
All numerical results were computed with \texttt{mpmath} at
50 decimal digits of precision, using \texttt{mpmath.zetazero}
for Riemann zero ordinates.
Matrix eigendecompositions and pseudoinverses were computed using
\texttt{NumPy} (\texttt{numpy.linalg}) with standard double-precision
floating point; the numerical rank of $\Tt$ is determined by a
threshold of $10^{-10}$ relative to the largest singular value.
Verification scripts are available at
\url{https://github.com/utehrani/analysislab-nt}.

\medskip
\emph{Not to be confused with.}
$\Ceta\neq\CT$: $\Ceta$ denotes the observed finite-grid ratio
$R_\eta(\kappa)=\lambda_{\max,\mathrm{ren}}/\pi(\kappa)\approx 0.39$
(no limit theorem is asserted), while $\CT$ governs the finite-grid
OLS fit $\mu_j\approx\CT(\kappa)/\gamma_{k(j)}$.
The inverse operator $W_1=\CT\cdot\Tt^+$ is \emph{not}
a Hilbert--P\'olya operator; its HP-correlation $r_2$
decays to $0.16$ with growing $\kappa$.
The renormalised asymmetry $\etaren$ uses weight $c_p^{\mathrm{ren}}=\sqrt{f_p}$;
Paper~3's $\etaorig$ uses $c_p^{\eta}=\sqrt{V_p(\tfrac12)}$:
these are numerically close but conceptually distinct.

\newpage

% -------------------------------------------------------
\section{Introduction}

\subsection{Background and Motivation}

This paper is the fourth in a series developing an operator-theoretic
framework for studying the distribution of the non-trivial zeros of
the Riemann zeta function. The operator-theoretic algebra in \S\S\,2--3
is non-circular and proof-complete; the spectral diagnostics in \S\,3
are numerical. \S\,4 contains proved estimates;
\S\,5 contains the numerical and conditional energy-asymmetry analysis.

\textbf{Paper~1}~\cite{Paper1} introduced the local curvature function and
established the divergence
$H_{\mathrm{local}}(\tfrac{1}{2},\kappa)\sim 2(\log\kappa)^2\to\infty$
at $\sigma=\tfrac{1}{2}$.

\textbf{Paper~2}~\cite{Paper2} connected this local curvature to the
Weil explicit functional, constructing admissible test functions
for which the curvature divergence becomes a structural feature of
the Weil energy.

\textbf{Paper~3}~\cite{Paper3} made the geometric structure explicit: two
finite-dimensional real Hilbert spaces $\Hstr$ and $\Hnull$ connected
by a linear map $\Phi$, with loop operator $T = \Phi^*\circ\Phi$ and
numerical energy asymmetry for the Paper~3 weight
$c_p^\eta=\sqrt{V_p(\tfrac12)}$ (denoted $\etaorig$ in Paper~3).
These results appear in \S\,1 for context; the formal core
\S\S\,2--5 depends only on the definitions in \S\,2.1.

\textbf{The present paper} introduces and analyzes $\Tt = \Phi\circ\Phi^*$,
the dual of $T$, acting on the zero space $\Hnull$. The operator
emerges canonically from the same $\Phi$ and requires no new
postulate. Several established programs seek a spectral interpretation
of the Riemann zeros; we position the present work relative to the
main ones.

A Hilbert--P\'olya operator $\mathcal{H}$ would satisfy
$\mathcal{H}\psi_k = \gamma_k\psi_k$; the zero ordinates appear
directly as eigenvalues. The Tehrani operator $\Tt$ has eigenvalues
$\mu_j\approx\CT(\kappa)/\gamma_{k(j)}$ (finite-grid OLS fit, numerically); the zero ordinates appear
as localization centers of eigenvectors, not as eigenvalues. We use the
standard formulation of the Hilbert--P\'olya setting (see e.g.\ the
survey by Bombieri~\cite{Bombieri} and the physics
review~\cite{SchumayerHutchinson11}) without adopting it as a postulate.
We call $\Tt$ a ``resonance operator'' to express eigenvector
concentration around arithmetic centers $\gamma_k$; this is
paper-specific terminology, not used in the scattering-theoretic
sense.\footnote{We call $\Tt=\Phi\Phi^*$ the \emph{resonance operator}
(Tehrani operator), following the notation of this series.
This should not be confused with resonance operators in the
Lax--Phillips scattering theory or with ``resonances'' in the
Sierra--Rodr\'iguez-Laguna model~\cite{SierraRodriguez11}, where
the term refers to poles of the resolvent.
Here $\Tt$ is a positive semidefinite Gram-type operator
on the finite-dimensional space $\Hnull$.
For analogous HP falsification results, see~\cite{EndresSteiner10}.}

Berry and Keating~\cite{BerryKeating} seek a self-adjoint operator
whose eigenvalues are $\gamma_k$ via $xp$-quantisation; the present
paper operates in the opposite direction ($\gamma_k$ as eigenvector
centers, not eigenvalues).
Burnol~\cite{BurnolForum04} constructs complete and minimal systems
in Sonin spaces indexed by primes and zeros via a conductor operator,
the closest published antecedent to the $\Hstr$--$\Hnull$ duality;
the present paper differs in using a finite-rank linear map $\Phi$
with explicit Gaussian phasors rather than an
infinite-dimensional Fourier-analytic conductor operator.
Connes~\cite{Connes} gives an adelic trace-formula interpretation;
no formal connection to $\Tt$ is established here;
for Weil positivity at the archimedean place, see
Connes--Consani~\cite{ConnesConsani21}.
Sierra and Townsend~\cite{SierraTownsend} provide a truncated
Hamiltonian realization connected to Connes' absorption spectrum;
$\Tt$ is not a Hamiltonian but a dual encoding operator.
The Bost--Connes system~\cite{BostConnes} operates in a
quantum-statistical framework with KMS states; the present
construction is purely finite-dimensional and linear-algebraic.
Most recently, Connes--Consani--Moscovici~\cite{CCM25} construct
a self-adjoint operator as a rank-one perturbation of the
scaling spectral triple, achieving numerical alignment with the
first fifty zero ordinates; our construction differs structurally:
$\Tt=\Phi\Phi^*$ is a Gram-type operator on $\Hnull$, indexed
by sin-phase prime phasors, without spectral postulates and
without rank-one perturbation structure.
The estimate, for each fixed $\gamma_k$,
$M_k(\kappa)\le(1/\sqrt{1+\gamma_k^2}+o_{\gamma_k}(1))\pi(\kappa)$
(no uniformity in $k$) is a standard prime exponential sum bound
proved by PNT/Abel-summation methods; see also Tur\'an's program~\cite{Turan}.
The closest historical precedent for the oscillatory-sum estimate
is Tur\'an's~\cite{Turan} program on approximative Dirichlet
polynomials; the novelty of our Proposition~\ref{thm:EXPLICIT} lies in how it
may enter a future analytic route to the energy-asymmetry problem.

The main conditional result is a proof that $\etaren > 0$ under
Assumption $(E_{\mathrm{rem}})$ --- that the cross/stream energy remainder
$|\Dcross+\Dstream|$ is bounded by $\rho\cdot\Dburstvar$ for some $\rho<1$.
Numerically $\rho_{\max}=0.583<1$ on the full tested grid
$\kappa\in\{23,53,101,199,503,1009\}$.

\subsection{Main Results}

We prove the algebraic spectral identity
$\sigma(T)\setminus\{0\}=\sigma(\Tt)\setminus\{0\}$ (\S\,2),
self-adjointness of $\Wone = \CT\cdot\Tt^+$ (\S\,3.4), the Abel
Summation Principle (Lemma~\ref{lem:M3}, \S\,4.3), and the Prime Exponential
Endpoint Estimate (\S\,4.4, Proposition~\ref{thm:EXPLICIT}) using
only the Prime Number Theorem. We establish numerically that $\Tt$
is not a Hilbert--P\'olya operator (\S\,2.4), that $\etaren > 0$
on the tested grid $\kappa\in\{23,53,101,199,503,1009\}$ (\S\,5.1), and that $\Dburstvar=3.105>0$
is $\kappa$-invariant by definition for $\kappa\ge 7$ (\S\,5.2);
$\Delta(\kappa)\ge\Dburstvar$ is observed on the tested grid
but not proved without additional assumptions.
Under Assumption $(E_{\mathrm{rem}})$,
$\etaren(\kappa)\ge(1-\rho)\,\Dburstvar/E_{\mathrm{str}}(\kappa)>0$ (\S\,5.3).

\subsection{Reader Contract}

This paper \emph{proves}, by classical operator theory and the Prime
Number Theorem, the spectral identity
$\sigma(T)\setminus\{0\}=\sigma(\Tt)\setminus\{0\}$,
self-adjointness of $\Wone$,
the Abel Summation Principle, and, for each fixed $\gamma_k$,
$M_k(\kappa)\le(1/\sqrt{1+\gamma_k^2}+o_{\gamma_k}(1))\pi(\kappa)$
with no uniformity in $k$ (Proposition~\ref{thm:EXPLICIT}).
These are unconditional.

This paper \emph{establishes numerically} that $\Tt$ is not a
Hilbert--P\'olya operator ($r_2$ falls to $0.16$ on the tested grid),
that $\etaren > 0$ on the tested grid
$\kappa\in\{23,53,101,199,503,1009\}$,
and that $\Dburstvar=3.105>0$ is
$\kappa$-invariant by definition for $\kappa\ge 7$.
These are verified on a finite grid and labelled
[numerical].

This paper \emph{proves conditionally/numerically}, under
Assumption $(E_{\mathrm{rem}})$ (\S\,5.3) and the finite numerical
verification $\Dburstvar>0$ at $(\varepsilon,N)=(0.05,100)$,
that $\etaren(\kappa)\ge(1-\rho)\,\Dburstvar/E_{\mathrm{str}}(\kappa)>0$.

This paper \emph{leaves open} seven problems: rank equality for $\Tt$
(Problem~6.1), the analytic proof of Assumption $(E_{\mathrm{rem}})$ (Problem~6.2),
the analytic proof of $\etaren > 0$ without assumptions (Problem~6.3),
a circularity-free HP spectral function $f(\Tt)$ (Problem~6.4), the
arithmetic origin of $\Ceta\approx 0.39$ (Problem~6.5), the
asymptotics of $\CT(\kappa)/\pi(\kappa)$ (Problem~6.6), and whether
$\etaren > 0$ follows from PNT alone (Problem~6.7).

% -------------------------------------------------------
\section{The Tehrani Operator}

\subsection{Recollection: $\Hstr$, $\Hnull$, $\Phi$, $T$}

We fix throughout:
\begin{itemize}
\item A prime cutoff $\kappa > 1$, with active prime set
  $S(\kappa) := \{p\text{ prime}: p\leq\kappa\}$
  and $\pi(\kappa) := |S(\kappa)|$.
\item The first $N$ positive imaginary parts of nontrivial zeros of
  $\zeta(s)$, ordered $0 < \gamma_1 \leq \gamma_2 \leq \cdots \leq \gamma_N$.
\item A Gaussian damping parameter $\varepsilon > 0$.
\end{itemize}

\begin{definition}[$\Hstr$ and $\Hnull$, from Paper~3]
\[
\Hstr := \ell^2(S(\kappa)),\quad \dim\Hstr = \pi(\kappa).
\qquad
\Hnull := \ell^2(\{\gamma_k\}_{k=1}^N),\quad \dim\Hnull = N.
\]
Both spaces carry standard inner products.
We work throughout with real vectors and real operators.
\end{definition}

\begin{definition}[Resonance vectors, from Paper~3]
For each prime $p\in S(\kappa)$, the \emph{resonance vector}
$a_p\in\Hnull$ is defined by
\[
a_p := \bigl(e^{-\varepsilon^2\gamma_k^2/2}\sin(\gamma_k\log p)\bigr)_{k=1}^N.
\]
Since $\Hnull$ is finite-dimensional ($\dim\Hnull=N$), $a_p\in\Hnull$
automatically; the Gaussian factor $e^{-\varepsilon^2\gamma_k^2/2}$ is
the damping used in the zero-coordinate representation.
\end{definition}

\begin{definition}[Linear map $\Phi$, from Paper~3]
\[
\Phi\colon\Hstr\to\Hnull,\quad
\Phi(e_p) = a_p,\quad
\Phi(c) = \sum_{p\leq\kappa}c_p a_p.
\]
\end{definition}

\begin{definition}[Loop operator $T$, from Paper~3]
\[
T := \Phi^*\circ\Phi\colon\Hstr\to\Hstr,\quad
T_{pq} = \langle a_p,a_q\rangle_{\Hnull} = G^{\mathrm{un}}_{pq}.
\]
$T$ is self-adjoint, positive semi-definite, and identical to the Gram
matrix $\Gun$ of the resonance vectors.
\end{definition}

\begin{definition}[Canonical weight vector, from Paper~2]\label{def:weight}
The canonical weight vector $c\in\Hstr$ used throughout this paper
has components
\[
  c_p \;:=\; \sqrt{\fp}, \qquad
  \fp \;:=\; \frac{4(\log p)^2(2p-1)}{p(p-1)^2},
\]
the Weil-renormalised weight from Paper~2.
It differs from Paper~3's $\eta$-framework weight
$c_p^\eta = \sqrt{V_p(\tfrac12)}$; see Paper~3, Notation and Conventions and \S\,4.
The leading asymptotic for large $p$ is
$c_p \sim 2\sqrt{2}\,\frac{\log p}{\sqrt{p}} + O\!\left(\frac{\log p}{p^{3/2}}\right)$,
since $f_p \sim 8(\log p)^2/p$.
The exact formula is used throughout this paper and in all verification scripts.
The local prime curvature $V_p(\sigma)$ of Paper~1 is related to $\fp$ by
$\fp = V_p(\tfrac12) - 4(\log p)^2/p$
(subtraction of the leading divergent term).

Reference values for the Paper~4 weight $c_p=\sqrt{f_p}$:
$c_p|_{p=2}\approx 1.698$, $c_p|_{p=3}\approx 1.418$,
$c_p|_{p=5}\approx 1.080$, $c_p|_{p=7}\approx 0.884$.
\end{definition}

\begin{definition}[Energy quantities, from Paper~3]
\[
\Estr := \sum_{p\leq\kappa}c_p^2\|a_p\|^2 = c^T \Da c,
\qquad
\Espec := \Bigl\|\sum_{p\leq\kappa}c_p a_p\Bigr\|^2 = \|\Phi c\|^2,
\]
where $\Da := \mathrm{diag}(\|a_p\|^2)_{p\in S(\kappa)}$.
The energy asymmetry is
\[
\Delta := \Estr - \Espec \geq 0\quad\text{(numerically)},
\qquad
\etaren := \frac{\Delta}{\Estr}
= 1 - \frac{c^T\Gun c}{\Estr}
= \eta[c^{\mathrm{ren}}].
\]
Here $\eta[c]$ is the asymmetry functional evaluated at weight $c$;
this paper uses $c=c^{\mathrm{ren}}:=(c_p^{\mathrm{ren}})_{p\le\kappa}$
with $c_p^{\mathrm{ren}}:=\sqrt{f_p}$.
\end{definition}

\subsection{Definition of the Tehrani Operator $\widetilde{T}$}

\begin{definition}[Tehrani operator]\label{def:Tt}
\[
\Tt := \Phi\circ\Phi^*\colon\Hnull\to\Hnull.
\]

\textbf{Matrix representation.}
In the standard basis $\{f_k\}_{k=1}^N$ of $\Hnull$:
\[
\Tt_{kl}
= \sum_{p\leq\kappa}(a_p)_k(a_p)_l
= \sum_{p\leq\kappa}
e^{-\varepsilon^2\gamma_k^2/2}\sin(\gamma_k\log p)
\cdot
e^{-\varepsilon^2\gamma_l^2/2}\sin(\gamma_l\log p).
\]
\end{definition}

\begin{proposition}[Properties of $\Tt$; proved]\label{prop:Tt-properties}
The operator $\Tt$ satisfies:
\begin{enumerate}[label=(\roman*)]
	\item \textbf{Self-adjointness:}
	$\Tt^* = (\Phi\Phi^*)^* = \Phi\Phi^* = \Tt$.
	\item \textbf{Positive semi-definiteness:}
	For all $u\in\Hnull$,
	$\langle\Tt u,u\rangle = \|\Phi^* u\|^2 \geq 0$.
	\item \textbf{Rank bound:}
	$\mathrm{rank}(\Tt) \leq \min\{N,\pi(\kappa)\}$.
\end{enumerate}
\emph{Proof.}
Parts (i) and (ii) are immediate from $\Tt=\Phi\Phi^*$ using standard
properties of adjoints.
Part (iii) follows from
$\mathrm{rank}(AB)\leq\min(\mathrm{rank}(A),\mathrm{rank}(B))$
applied to $\Tt=\Phi\Phi^*$, giving
$\mathrm{rank}(\Tt)\leq\min(N,\pi(\kappa))$
(the bound from the singular values of $\Phi\in\mathbb{R}^{N\times m}$).
The reverse inequality $\mathrm{rank}(\Tt)\geq\pi(\kappa)$ would require
linear independence of $\{a_p\}_{p\leq\kappa}$ in $\Hnull$; this is
not proved analytically and is listed in \S\ref{sec:open-problems}.
$\square$
\end{proposition}

\begin{remark}[Non-circularity of $\Tt$]
$\Tt$ is built directly from prime resonance vectors $\{a_p\}$.
The zero ordinates $\gamma_k$ are numerical inputs to the resonance
vectors and hence to $\Tt$.
Non-circularity means: no RH, GUE, or HP assumption is used; the
$\gamma_k$ are external inputs, not conclusions of the model.
\end{remark}

\emph{Numerical evidence.} At benchmark parameters
$(\kappa,\varepsilon,N)=(53,0.05,100)$ the spectrum of $\Tt$ exhibits
a sharp gap between the $\pi(\kappa)=16$ largest eigenvalues
($\mu_{16}\approx 4.94\cdot 10^{-6}$) and the remaining ones
($\mu_{17}\approx 1.81\cdot 10^{-18}$), giving the ratio
$\mu_{17}/\mu_{16}\approx 3.7\cdot 10^{-13}$. The numerical rank
therefore equals $\pi(\kappa)$ at machine precision, consistent
with linear independence of $\{a_p\}$ and hence with
$\mathrm{rank}(\Tt)=\pi(\kappa)$, but the analytic statement remains
an open problem.

\subsection{Algebraic Spectral Identity}

\begin{theorem}[Spectral identity; proved]\label{thm:spectral-identity}
\[
\sigma(T)\setminus\{0\} = \sigma(\Tt)\setminus\{0\}.
\]
\emph{Proof.}
Let $\Phi$ be represented by an $N\times m$ real matrix
($m=\pi(\kappa)$). The nonzero eigenvalues of $T=\Phi^*\Phi$
and $\Tt=\Phi\Phi^*$, counted with algebraic multiplicity,
are the squared nonzero singular values of $\Phi$.
Hence $\sigma(T)\setminus\{0\}=\sigma(\Tt)\setminus\{0\}$.
The zero-eigenvalue multiplicities differ by the dimension gap:
if $N\ge m$, then $\Tt$ has $N-m$ additional zero eigenvalues;
if $m>N$, then $T$ has $m-N$ additional zero eigenvalues.
At the benchmark $(\kappa,N)=(53,100)$, this gives
$N-m=100-16=84$ additional zero eigenvalues for $\Tt$. $\square$
\end{theorem}

\begin{remark}[Numerical implementation check]
At benchmark parameters $(\kappa,\varepsilon,N)=(53,0.05,100)$:
\[
\max_{j=1}^{\pi(\kappa)}|\lambda_j(T) - \mu_j(\Tt)| < 10^{-15}
\quad\text{(machine precision)}.
\]
This verifies the algebraic identity numerically.
\end{remark}

\begin{remark}
The identity $\lambda_j(T)=\mu_j(\Tt)$ is an \emph{algebraic} theorem
--- it holds in exact arithmetic for any choice of $\kappa$, $\varepsilon$,
and $N$.
The eigenvalues $\lambda_j = \mu_j$ are not the zero ordinates $\gamma_k$;
they are constructive spectral quantities arising from $\Phi$.
\end{remark}

\subsection{$\Tt$ is NOT a Hilbert--P\'olya Operator}

This section records a negative result.

\begin{proposition}[Hilbert--P\'olya diagnostic failure; numerical]\label{prop:not-HP}
On the tested finite-cutoff grid, $\Tt$ and $\Wone=\CT\Tt^+$
give no numerical support for an asymptotic Hilbert--P\'olya
interpretation:
\begin{enumerate}[label=(\alph*)]
\item \emph{[Numerical]} The eigenvalues of $\Tt$ satisfy
  $\mu_j\approx\CT(\kappa)/\gamma_{k(j)}$ (finite-grid OLS fit) --- they are \emph{inversely} proportional
  to zero ordinates (\S\,3.1).
\item \emph{[Numerical]} The zero ordinates $\gamma_k$ appear as
  \emph{localization centers} of eigenvectors $v_j$ (\S\,3.2),
  not as eigenvalues.
\item \emph{[Numerical]} The inverse operator $\Wone := \CT\cdot\Tt^+$
  has eigenvalues $\omega_j = \CT/\mu_j \approx \gamma_{k(j)}$
  for small $j$, but the correlation
  $r_2 := \mathrm{corr}(\omega_j, \gamma_{k(j)})$
  falls from $0.502$ (at $\kappa=23$) to $0.160$ (at $\kappa=503$).
  On the tested grid, $\Wone$ fails the HP-correlation diagnostic;
these data give no numerical support for $\Wone$ as an
asymptotic Hilbert--P\'olya operator.
\end{enumerate}
\emph{The fundamental distinction:}
\[
\text{HP-operator: eigenvalues} = \gamma_k.
\qquad
\Tt\text{: eigenvector localization centers} = \gamma_k.
\]
$\Tt$ is a \textbf{resonance operator}: it encodes the $\gamma_k$
as mode centers, not as spectral values.

\emph{Proof of (a)(b)(c).}
Numerical; see \S\,3.1, \S\,3.2, \S\,3.4 for detailed data. $\square$
\end{proposition}

\begin{remark}[Numerical obstruction pattern]
Gaussian damping $e^{-\varepsilon^2\gamma_k^2}$ is a plausible
explanation for the finite-grid failure of the HP-correlation
diagnostic. For $j=16$ at $\kappa=53$, the dominant zero is $\gamma_1$,
and the HP-correlation $r_2$ falls to $0.16$ on the tested grid.
This finite-grid pattern is consistent with the observed
failure of $\Wone$ as an asymptotic HP-candidate;
no structural obstruction theorem is proved here.
\end{remark}

% -------------------------------------------------------
\section{Spectral Analysis of $\widetilde{T}$}

\subsection{Eigenvalue Structure: $\mu_j \approx \CT(\kappa)/\gamma_{k(j)}$ (finite-grid OLS)}

\emph{Numerical evidence} (Eigenvalue--Zero Ordinate Correlation).

Let $\kdom(j) := \arg\max_k|(v_j)_k|$ be the dominant zero index
for eigenvector $v_j$. Then:

\begin{center}
\begin{tabular}{ccccc}
\toprule
$j$ & $\mu_j(\Tt)$ & $\kdom$ & $\gamma_{\kdom}$ & Localization \\
\midrule
1 & 4.219803 & 1 & 14.1347 & 0.985 \\
2 & 2.565694 & 3 & 25.0109 & 0.741 \\
3 & 1.744649 & 2 & 21.0220 & 0.745 \\
4 & 0.986492 & 4 & 30.4249 & 0.854 \\
5 & 0.314199 & 5 & 32.9351 & 0.794 \\
6 & 0.217900 & 6 & 37.5862 & 0.885 \\
\bottomrule
\end{tabular}
\end{center}

The OLS correlation between $\mu_j$ and $1/\gamma_{k(j)}$ is
\[
r_1 = \mathrm{corr}\bigl(\mu_j,\,1/\gamma_{k(j)}\bigr) = 0.950
\quad(\kappa=53).
\]
For $\kappa=503$: $r_1 = 0.859$ (weakening as $\kappa$ grows,
consistent with Gaussian damping).

\begin{remark}
The dominant index satisfies $\kdom(1) = 1$ for all tested
$\kappa\in\{23,53,101,199,503,1009\}$:
the first eigenvector of $\Tt$ is always associated with
$\gamma_1 = 14.1347$.
\end{remark}

\begin{remark}[Cauchy--Schwarz warning]
Note that $r_1 = 0.950$ coexists with $r_2 = 0.502$ (\S\,3.4).
These measure different things: $r_1$ tests monotone ordering;
$r_2$ tests value-level approximation.
The distinction between ordering and value accuracy is crucial for
the HP assessment in \S\,2.4.
\end{remark}

\subsection{Eigenvector Localization}

\emph{Numerical evidence} (Eigenvector Localization).

For each eigenvector $v_j$ of $\Tt$ (normalized: $\|v_j\|=1$),
the localization is measured by $L_j := |(v_j)_{\kdom(j)}|$.
Observed range: $L_j\in[0.45,0.99]$ for $j=1,\ldots,14$.

For comparison, a uniformly random unit vector in $\mathbb{R}^{100}$
has typical entry magnitude $1/\sqrt{100} = 0.10$.
The observed localization is $4.5$--$9.9$ times the random baseline.
The localization structure suggests that $\Tt$ encodes
the ordering of zero ordinates in a distributional sense:
each mode $v_j$ is concentrated near the zero $\gamma_{k(j)}$
in the zero-space $\Hnull$.
This is the precise content of calling $\Tt$ a resonance operator.

\begin{assumption}[Nonzero resonance norms]\label{ass:Da-pos}
For the finite parameter sets under consideration,
$\|a_p\|_{\Hnull}>0$ for every active prime $p\le\kappa$.
Equivalently, $\Da$ is positive diagonal, and $\Da^{-1/2}$
is well-defined. This is verified numerically for all
parameter sets used below.
\end{assumption}

\subsection{The Two Arithmetic Constants $\Ceta$ and $\CT$}

\begin{definition}[Arithmetic constants]\label{def:constants}
\[
R_\eta(\kappa) := \frac{\lambda_{\max}(\Tren)}{\pi(\kappa)},
\qquad
\Ceta \approx 0.39\quad\text{(observed finite-grid value; no limit theorem)},
\]
where $\Tren := \Da^{-1/2}T\Da^{-1/2}$.
\[
\CT(\kappa) := \text{OLS-slope of } \mu_j\sim\CT/\gamma_{k(j)},
\qquad \CT/\pi(\kappa) \in \{2.49,\,2.21,\,2.17,\,1.94,\,1.89\}
\]
for $\kappa\in\{23,53,101,199,503\}$ (weakly decreasing).

\[
\Ceta \neq \CT.
\]
\emph{These are two distinct arithmetic quantities; always use subscripts.}
\end{definition}

More explicitly, let $J(\kappa)=\{1,\ldots,r(\kappa)\}$ where
$r(\kappa)=\operatorname{rank}_{\mathrm{num}}(\Tt)$, and let
$\kdom(j)=\arg\max_k|(v_j)_k|$. Then:
\[
\CT(\kappa) := \arg\min_{C\in\mathbb{R}}
\sum_{j\in J(\kappa)}
\Bigl(\mu_j - \frac{C}{\gamma_{\kdom(j)}}\Bigr)^2.
\]
No intercept term is used.

The ratio $R_\eta(\kappa)$ is numerically stable at $\Ceta\approx 0.39$
across $\kappa\in\{23,\ldots,1009\}$ (see table; $\sigma^2 < 0.0003$).
Whether $R_\eta(\kappa)$ converges to a limit is analytically open;
we write $\Ceta$ as a name for the observed value.

\emph{Numerical evidence} ($\lambda_{\max,\mathrm{ren}}$ growth).

\begin{center}
\begin{tabular}{cccc}
\toprule
$\kappa$ & $\pi(\kappa)$ & $\lambda_{\max,\mathrm{ren}}$ & Ratio \\
\midrule
23  & 9   & 3.59  & 0.399 \\
53  & 16  & 6.38  & 0.399 \\
101 & 26  & 10.35 & 0.398 \\
199 & 46  & 17.76 & 0.386 \\
503 & 96  & 37.38 & 0.389 \\
1009 & 169 & 65.16 & 0.386 \\
\bottomrule
\end{tabular}
\end{center}

Variance of the ratio: $\sigma^2 < 0.0003$ --- extremely stable.
The universal bound $\lambda_{\max,\mathrm{ren}} < 1$
(an earlier conjecture in this series, now falsified) is numerically falsified:
$\lambda_{\max,\mathrm{ren}} > 1$ already on the tested grid.
The data are compatible with growth proportional to $\pi(\kappa)$,
with observed ratio $R_\eta(\kappa)\approx 0.39$ (see Definition~\ref{def:constants});
convergence of this ratio remains analytically open.

\begin{remark}[Analytical status of $\Ceta$]
The observed ratio $R_\eta(\kappa)\approx 0.39$ (denoted $\Ceta$)
is not analytically explained.
A candidate formula $\Ceta = 1/I_N(\varepsilon)$ is \textbf{falsified}:
it gives $1/I_{100}(0.05)\approx 0.73\neq 0.39$.
The correct representation involves a harmonic mean
(analytically open; see Open Problem~\ref{op:ceta}).
\end{remark}

\subsection{The Inverse Operator $\Wone = \CT\cdot\Tt^+$}

\begin{definition}[Inverse operator]
$\Wone := \CT\cdot\Tt^+$,
where $\Tt^+$ denotes the Moore--Penrose pseudoinverse of $\Tt$
and $\CT$ is the OLS constant from Definition~\ref{def:constants}.
\end{definition}

\begin{proposition}[Self-adjointness of $\Wone$; proved]\label{prop:W1-sa}
$\Wone$ is self-adjoint.

\emph{Proof.}
$\Tt$ is self-adjoint (Proposition~\ref{prop:Tt-properties}(i)).
For any self-adjoint operator, the Moore--Penrose pseudoinverse is also
self-adjoint (standard result in linear algebra).
$\CT$ is a real scalar.
Therefore $\Wone = \CT\cdot\Tt^+$ is self-adjoint. $\square$
\end{proposition}

Numerically: $\|\Wone - \Wone^T\|_\infty = 3.1\times 10^{-4}$,
$\|\Wone - \Wone^T\|_\infty/\|\Wone\|_\infty \approx 6.7\times 10^{-11}$
(rounding from pseudoinversion of a moderately conditioned matrix).

\begin{proposition}[Non-circularity status of $\Wone$; proved]\label{prop:W1-circ}
	$\Wone = \CT\cdot\Tt^+$ does not assume RH or the Hilbert--P\'olya postulate.
	However, its normalization constant $\CT$ is the OLS slope in
	$\mu_j\approx\CT(\kappa)/\gamma_{k(j)}$ (finite-grid OLS fit),
	which is calibrated against the zero ordinates.
	$\Wone$ should therefore not be presented as fully input-free.
	$\Tt$ is explicit and deterministic; the $\gamma_k$ are external
	numerical inputs, not conclusions of the model.
	\emph{Proof.}
	$\Tt$ is constructed from prime resonance vectors
	(Definition~\ref{def:Tt}), using $\gamma_k$ as numerical inputs.
	$\Tt^+$ is the Moore--Penrose pseudoinverse --- no postulate.
	The scalar $\CT$ is computed via OLS regression of $\mu_j$ against
	$1/\gamma_{k(j)}$: this uses $\gamma_k$ as calibration values, but
	does not assume RH or HP. $\square$
\end{proposition}

\emph{Numerical evidence} ($\Wone$ HP-correlation failure).

\begin{center}
\begin{tabular}{ccc}
\toprule
$\kappa$ & $\pi(\kappa)$ & $r_2 = \mathrm{corr}(\omega_j,\gamma_{k(j)})$ \\
\midrule
23  & 9  & 0.502 \\
53  & 16 & 0.436 \\
101 & 26 & 0.321 \\
199 & 46 & 0.226 \\
503 & 96 & 0.160 \\
\bottomrule
\end{tabular}
\end{center}

The decreasing trend indicates, on the tested grid,
that $\Wone$ is not supported as an asymptotic
HP-operator.
The observed pattern is consistent with Gaussian damping:
$e^{-\varepsilon^2\gamma_k^2/2}$ in $\Phi$ is a plausible explanation
for the rapid decay of $\mu_j$, with a fitted
power-law exponent $\alpha\approx -4.88$.
This decay is much steeper than the approximate $1/\gamma$ behaviour
that would be needed for the pseudoinverse eigenvalues
$\omega_j=\CT/\mu_j$ to track the zero ordinates.
Pseudoinversion therefore strongly amplifies tail modes.

\textbf{Summary for \S\,3.4:}
$\Wone$ is self-adjoint (proved).
Its construction does not assume RH or the HP postulate, but $\CT$
is calibrated via OLS on the zero ordinates --- $\Wone$ is therefore
not fully input-free (Proposition~\ref{prop:W1-circ}).
$\Wone$ is the natural HP-candidate but fails numerically as
$\kappa$ grows.
These diagnostics give no numerical support for
this particular pseudoinverse construction as an
asymptotic Hilbert--P\'olya candidate.
This does not preclude other spectral functions $f(\Tt)$ from being valid
HP-candidates; see \S\,6.

% -------------------------------------------------------
\section{Monotonicity and Prime Exponential Sums}

The following estimates record ingredients for a possible
analytic route toward bounding $\etaren$ from below,
based on the interplay of three structural features:
the monotone weight sequence $(c_p)$,
the oscillatory mode structure $(\sin(\gamma_k\log p))$,
and Abel summation.

\subsection{Monotone weight structure}

\begin{lemma}[Monotone weight structure; proved for $p\geq 11$]\label{lem:M1}

The canonical weight $c_p = \sqrt{f_p}$ is strictly monotone
decreasing in $p$ for all primes $p\geq 11$.

\emph{Proof.}
The exact expression for $c_p$ is
\[
c_p = \frac{2\log p\,\sqrt{2p-1}}{\sqrt{p}\,(p-1)}.
\]
Taking the logarithmic derivative:
\[
\frac{d}{dp}\log c_p
= \frac{1}{p\log p} + \frac{1}{2p-1} - \frac{1}{2p} - \frac{1}{p-1}.
\]
For $p\ge 11$, we have $\log p > 2$, hence $1/(p\log p) < 1/(2p)$.
Therefore:
\[
\frac{d}{dp}\log c_p
< \frac{1}{2p} + \frac{1}{2p-1} - \frac{1}{2p} - \frac{1}{p-1}
= \frac{1}{2p-1} - \frac{1}{p-1} < 0,
\]
since $2p-1 > p-1 > 0$. Hence $c_p$ is strictly decreasing
in $p$ for all $p\ge 11$. $\square$
\end{lemma}

\subsection{Oscillatory mode structure}

\begin{lemma}[Oscillatory mode intervals; proved]\label{lem:M2a}
For fixed $\gamma_k > 0$, the function $p\mapsto\sin(\gamma_k\log p)$
has zeros at $p = e^{m\pi/\gamma_k}$ for integers $m\ge 0$,
and its sign is constant on each open interval
$(e^{m\pi/\gamma_k},\, e^{(m+1)\pi/\gamma_k})$,
alternating between adjacent intervals.
Restricting to primes $p\ge 2$ selects the relevant subcollection.

\emph{Proof.}
Direct from $\sin(x)=0$ at $x=m\pi$: the argument
$\gamma_k\log p$ crosses $m\pi$ at $p = e^{m\pi/\gamma_k}$. $\square$
\end{lemma}

\begin{remark}[Prime occupation; open]
For fixed $\gamma_k$, the PNT suggests eventual prime occupation
of the intervals $I_m$ as $\kappa\to\infty$; a quantitative
finite-cutoff version with the required uniformity is not proved here.
\end{remark}

\subsection{Abel Summation Principle}

\begin{lemma}[Abel summation principle; proved]\label{lem:M3}
Let $P\geq 1$.
Let $(f_1,\ldots,f_P)$ be monotone non-increasing and non-negative.
Let $(g_1,\ldots,g_P)$ be any sequence with partial sums bounded by $M$:
$|G_m| := |\sum_{i=1}^m g_i|\leq M$ for all $m$.
Then:
\[
\Bigl|\sum_{i=1}^P f_i g_i\Bigr| \leq M\cdot f_1.
\]
\emph{Proof.}
Apply the Abel partial summation formula:
$\sum_{i=1}^P f_i g_i = f_P G_P - \sum_{i=1}^{P-1}(f_{i+1}-f_i)G_i$.
Taking absolute values:
\begin{align*}
\Bigl|\sum_{i=1}^P f_i g_i\Bigr|
&\leq |f_P|\cdot|G_P| + \sum_{i=1}^{P-1}|f_{i+1}-f_i|\cdot|G_i| \\
&\leq Mf_P + M\sum_{i=1}^{P-1}(f_i-f_{i+1})
 \quad\text{[$f$ decreasing]} \\
&= Mf_P + M(f_1-f_P) = Mf_1. \quad\square
\end{align*}
\end{lemma}

\begin{remark}[Application context for Paper~4]
Lemma~\ref{lem:M1} establishes monotonicity of the unnormalised
weights $c_{p_i}$.
In the Rayleigh quotient identity of \S\,5.1, the relevant
normalised vector is $\tilde{c} = \Da^{1/2}c/\sqrt{\Estr}$,
whose components satisfy $(\tilde{c})_p \propto c_p\|a_p\|$.
Monotonicity of $c_p$ therefore does not by itself imply
monotonicity of the weights entering the spectral projections,
since $\|a_p\|$ depends on the zero-sum profile and may
disturb the ordering.
Applying Lemma~\ref{lem:M3} to spectral projections therefore
requires either a separate monotonicity or bounded-variation
estimate for $c_p\|a_p\|$, or a different transfer lemma from
prime-phasor partial sums to the normalised spectral projections
$|\langle u_j,\tilde{c}\rangle|$.
This is the missing transfer step noted in Open Problem~\ref{op:pos}.
\end{remark}

\subsection{Prime Exponential Endpoint Estimate}

\begin{proposition}[Prime exponential endpoint estimate; proved]\label{thm:EXPLICIT}
Let $t\ne 0$ be fixed. As $x\to\infty$:
\[
\sum_{p\le x}p^{it}
= \frac{x^{1+it}}{(1+it)\log x}
  + O_t\!\left(\frac{x}{\log^2 x}\right).
\]
In particular, for the sine sum:
\[
M_k(x) := \max_{y\le x}\Bigl|\sum_{p\le y}\sin(\gamma_k\log p)\Bigr|
\le \left(\frac{1}{\sqrt{1+\gamma_k^2}}+o_{\gamma_k}(1)\right)\pi(x).
\]

\emph{Proof.}
Write $\vartheta(x)=\sum_{p\le x}\log p$ (Chebyshev's function).
By Abel partial summation, $\sum_{p\le x}p^{it}
= \int_2^x u^{it}\,\frac{d\vartheta(u)}{\log u}$.
The PNT gives $\vartheta(x)=x+o(x)$, and with a quantitative
error term $\vartheta(x)=x+O(xe^{-c\sqrt{\log x}})$
(Hadamard/de~la~Vall\'ee-Poussin); see \cite{IwaniecKowalski}, \S\,6.
Writing $E(u)=\vartheta(u)-u$, Stieltjes integration by parts gives
\begin{align*}
\int_2^x F(u)\,dE(u)
&= F(x)E(x) - F(2)E(2)
  - \int_2^x E(u)F'(u)\,du,
\end{align*}
where $F(u)=u^{it}/\log u$ and
$F'(u)=\bigl(it\log u - 1\bigr)u^{it-1}/(\log u)^2
       = O_t\bigl(1/(u\log u)\bigr)$.
The de~la~Vall\'ee-Poussin bound
$E(u)=O\bigl(ue^{-c\sqrt{\log u}}\bigr)$ implies
$E(x)F(x)=O_t\bigl(xe^{-c\sqrt{\log x}}\bigr)$
and
$\int_2^x |E(u)F'(u)|\,du
 = O_t\!\left(\int_2^x e^{-c\sqrt{\log u}}\,du\right)
 = O_t\!\left(xe^{-c'\sqrt{\log x}}\right)$.
Hence $\int_2^x F(u)\,dE(u)=O_t(x/\log^2 x)$,
and the main term comes from
$\int_2^x u^{it}/\log u\,du$,
yielding the displayed asymptotic.
Integrating by parts yields the first display:
\[
\int_2^x \frac{u^{it}}{\log u}\,du
= \frac{x^{1+it}}{(1+it)\log x}
  + O_t\!\left(\frac{x}{\log^2 x}\right),
\]
where the de~la~Vall\'ee-Poussin error is absorbed into the
same $O_t$-term for fixed $t$.
The estimate holds for fixed $t$; no uniformity in $t$ or zero index $k$
is asserted.
The sine sum bound follows by taking imaginary parts:
\[
\Bigl|\sum_{p\le x}\sin(t\log p)\Bigr|
= \Bigl|\operatorname{Im}\!\sum_{p\le x}p^{it}\Bigr|
\le \frac{x}{(1+t^2)^{1/2}\log x} + O_t\!\!\left(\frac{x}{\log^2 x}\right)
= \left(\frac{1}{\sqrt{1+t^2}}+o_t(1)\right)\pi(x). \quad\square
\]
The same estimate applies with $y$ in place of $x$ for any $y$ in the range.
For fixed $t$, the constant in the $O_t(\cdot)$-term is uniform
for all $y\ge y_0(t)$. Since $y/\log y$ and $y/\log^2 y$ are
increasing beyond a fixed threshold, the same bound holds after
taking the supremum over $y_0(t)\le y\le x$.
The finitely many values $2\le y < y_0(t)$ contribute $O_t(1)$,
hence $o_t(\pi(x))$ as $x\to\infty$; they do not affect the
maximum asymptotic.
\end{proposition}

\begin{remark}[Methodological note]
The estimate, for each fixed $\gamma_k$,
$M_k(\kappa)\le(1/\sqrt{1+\gamma_k^2}+o_{\gamma_k}(1))\pi(\kappa)$
(no uniformity in $k$) is a standard consequence of the Prime Number Theorem; the
analytic method has precedents in Tur\'an's program on
approximative Dirichlet polynomials \cite{Turan}.
What is specific to this paper is how this bound may enter
a future analytic route to the energy-asymmetry problem.
\end{remark}

\begin{remark}[What Proposition~\ref{thm:EXPLICIT} does NOT prove]
Proposition~\ref{thm:EXPLICIT} does not prove
$M_k(\kappa) = o(\pi(\kappa))$ for \emph{fixed} $k$ as $\kappa\to\infty$.
This stronger statement is related to the pair-cancellation
hypothesis $(E_{\mathrm{pair}})$ stated in \S\,5.3.
\end{remark}

% -------------------------------------------------------
\section{Results on $\etaren > 0$}

\subsection{Numerical Confirmation}

\emph{Numerical evidence} ($\etaren > 0$ on the tested grid $\kappa\in\{23,53,101,199,503,1009\}$).

Let $\tilde{c} = \Da^{1/2}c/\sqrt{\Estr}$ be the normalized canonical
weight vector. The Rayleigh quotient identity
\[
\langle\Tren\,\tilde{c},\,\tilde{c}\rangle = 1 - \etaren
\]
holds to machine precision ($< 2\times 10^{-16}$),
where $\Tren := \Da^{-1/2}T\Da^{-1/2}$.

\begin{center}
\begin{tabular}{ccccc}
\toprule
$\kappa$ & $\pi(\kappa)$ & $\langle\Tren\tilde{c},\tilde{c}\rangle$
  & $< 1$? & $\etaren$ \\
\midrule
23   & 9   & 0.18849 & $\checkmark$ & 0.81151 \\
53   & 16  & 0.33073 & $\checkmark$ & 0.66927 \\
101  & 26  & 0.22803 & $\checkmark$ & 0.77197 \\
199  & 46  & 0.21740 & $\checkmark$ & 0.78260 \\
503  & 96  & 0.21788 & $\checkmark$ & 0.78212 \\
1009 & 169 & 0.18002 & $\checkmark$ & 0.81998 \\
\bottomrule
\end{tabular}
\end{center}

$\etaren > 0$ on the tested grid $\kappa\in\{23,53,101,199,503,1009\}$ $\checkmark$.

\textbf{Collective mechanism.}
Writing $\Tren u_j = \lambda_j u_j$, define the mode contribution
\[
q_j := \lambda_j\,|\langle u_j,\tilde{c}\rangle|^2,
\]
so that $\langle\Tren\tilde{c},\tilde{c}\rangle = \sum_j q_j$.
At $\kappa=53$: $q_4=0.193$, $q_1=0.040$.
The dominant contribution comes from mode $j=4$,
not from the dominant eigenmode $j=1$.
Approximately $92\%$ cancellation occurs at $j=1$ ($\kappa=53$).
The effect is distributed across $\sim 8$ eigenmodes ---
$\etaren > 0$ is a \textbf{multi-mode Rayleigh-quotient contribution}, not a simple
$u_1$-orthogonality principle.

\subsection{$\Delta$-Decomposition and Burst Stability}

\begin{lemma}[Exact energy-gap decomposition; proved]\label{lem:Delta-decomp}
Let $S(r):=\sum_{k=1}^N e^{-\varepsilon^2\gamma_k^2}\cos(\gamma_k\log r)$.
Then
\[
\Delta = \Dburstvar + \Dcross + \Dstream,
\]
where
\begin{align*}
\Dburstvar &=
  \sum_{\substack{p<q\le7\\p,q\text{ prime}}}
  c_p\,c_q\,[S(pq)-S(p/q)],\\
\Dcross &=
  \sum_{p\le7<q\le\kappa}
  c_p\,c_q\,[S(pq)-S(p/q)],\\
\Dstream &=
  \sum_{7<p<q\le\kappa}
  c_p\,c_q\,[S(pq)-S(p/q)].
\end{align*}
\end{lemma}

\begin{proof}
By definition,
\[
E_{\mathrm{spec}}
=\Bigl\|\sum_{p\le\kappa}c_p a_p\Bigr\|^2
=\sum_{p,q\le\kappa}c_p c_q\langle a_p,a_q\rangle.
\]
Since $E_{\mathrm{str}}=\sum_{p\le\kappa}c_p^2\|a_p\|^2$, the diagonal terms ($p=q$)
cancel in $\Delta=E_{\mathrm{str}}-E_{\mathrm{spec}}$, giving
\[
\Delta = -2\sum_{p<q\le\kappa}c_p c_q\langle a_p,a_q\rangle.
\]
Using $\cos(A+B)-\cos(A-B)=-2\sin A\sin B$ with
$A=\gamma_k\log p$, $B=\gamma_k\log q$:
\[
-2\langle a_p,a_q\rangle
= \sum_{k=1}^N e^{-\varepsilon^2\gamma_k^2}
  \bigl[\cos\bigl(\gamma_k\log(pq)\bigr)-\cos\bigl(\gamma_k\log(p/q)\bigr)\bigr]
= S(pq)-S(p/q).
\]
Splitting the double sum by ranges of $p,q$ gives the asserted decomposition. $\square$
\end{proof}

\emph{Numerical evidence} ($\Dburstvar$ stability).

\begin{center}
\begin{tabular}{lcccc}
\toprule
Pair & $S(pq)$ & $S(p/q)$ & $c_p c_q$ & Contribution \\
\midrule
$(2,3)$ & 0.9536 & 0.2653 & $\approx 2.408$ & $+1.657$ \\
$(2,5)$ & 0.4110 & 0.6555 & $\approx 1.833$ & $-0.448$ \\
$(2,7)$ & 0.5794 & 0.5812 & $\approx 1.501$ & $-0.003$ \\
$(3,5)$ & 0.9440 & 0.3686 & $\approx 1.531$ & $+0.881$ \\
$(3,7)$ & 0.7193 & 0.5571 & $\approx 1.254$ & $+0.203$ \\
$(5,7)$ & 0.9668 & 0.1135 & $\approx 0.954$ & $+0.814$ \\
\midrule
\textbf{Total} & & & & $\mathbf{\approx +3.11}$ \\
\bottomrule
\end{tabular}
\end{center}

At fixed $(\varepsilon,N)$, $\Dburstvar$ is independent of $\kappa$
for all $\kappa\ge 7$ by definition (the sum is over the fixed pairs
$p<q\le 7$). Numerically, at $(\varepsilon,N)=(0.05,100)$:
$\Dburstvar = 3.105\ldots > 0$ (Lemma~\ref{lem:burst-pos}).
The sufficient condition for $\Delta > 0$ is:
\[
|\Dcross + \Dstream| < \Dburstvar = 3.105.
\]
This is numerically satisfied on the tested grid
$\kappa\in\{23,53,101,199,503,1009\}$;
analytic proof beyond this grid is OPEN.

\subsection{Conditional Result under $(E_{\mathrm{rem}})$}

\begin{assumption}[$(E_{\mathrm{rem}})$; remainder control]\label{ass:Erem}
At fixed reference parameters $(\varepsilon,N)=(0.05,100)$
and with the Paper-4 weight $c_p^{\mathrm{ren}}=\sqrt{f_p}$,
there exist $\rho < 1$ and $\kappa_0\ge 7$ such that
for all $\kappa\ge\kappa_0$:
\[
\left|\Dcross(\kappa) + \Dstream(\kappa)\right|
\le \rho\cdot\Dburstvar.
\]
\end{assumption}

\begin{lemma}[Burst positivity; numerical]\label{lem:burst-pos}
At benchmark parameters $(\varepsilon,N)=(0.05,100)$,
the finite six-pair burst sum satisfies
\[
\Dburstvar
= \sum_{\substack{p<q\le7\\p,q\;\mathrm{prime}}}
  c_p^{\mathrm{ren}}\,c_q^{\mathrm{ren}}\,
  \bigl(S(pq)-S(p/q)\bigr)
= 3.105 \in [3.09,\,3.12].
\]
In particular $\Dburstvar > 0$.
Numerical verification code is available at
\url{https://github.com/utehrani/analysislab-nt}.
\end{lemma}

\begin{theorem}[Conditional/numerical positivity from remainder control]\label{thm:T5}
At the reference parameters
$(\sigma_0,\varepsilon,N)=(\tfrac12,0.05,100)$,
under the numerical bound $\Dburstvar>0$ (Lemma~\ref{lem:burst-pos}),
the exact decomposition of Lemma~\ref{lem:Delta-decomp},
Assumption~$(E_{\mathrm{rem}})$,
and $E_{\mathrm{str}}(\kappa)>0$
(Assumption~\ref{ass:Da-pos} and $c_p^{\mathrm{ren}}>0$),
for all $\kappa\ge\kappa_0$:
\[
\Delta(\kappa)
\ge (1-\rho)\,\Dburstvar > 0,
\]
and consequently
\[
\etaren(\kappa)
= \frac{\Delta(\kappa)}{E_{\mathrm{str}}(\kappa)}
\ge \frac{(1-\rho)\,\Dburstvar}{E_{\mathrm{str}}(\kappa)} > 0.
\]

\emph{Proof.}
By the triangle inequality and $(E_{\mathrm{rem}})$:
\[
\Delta(\kappa)
= \Dburstvar
  + \Dcross(\kappa)
  + \Dstream(\kappa)
\ge \Dburstvar
  - |\Dcross(\kappa)+\Dstream(\kappa)|
\ge (1-\rho)\,\Dburstvar > 0.
\]
By Assumption~\ref{ass:Da-pos} and $c_p^{\mathrm{ren}}>0$,
$E_{\mathrm{str}}(\kappa)>0$; hence division by $E_{\mathrm{str}}(\kappa)$ gives
$\etaren(\kappa) \ge (1-\rho)\,\Dburstvar/E_{\mathrm{str}}(\kappa) > 0$.
$\square$
\end{theorem}

\begin{remark}[Numerical verification of $(E_{\mathrm{rem}})$]
On the tested grid $\kappa\in\{23,53,101,199,503,1009\}$
with $\varepsilon=0.05$, $N=100$:
\begin{center}
\begin{tabular}{ccccc}
\toprule
$\kappa$ & $\Delta(\kappa)$ & $\Dburstvar$
& $|\Dcross+\Dstream|$ & $\rho(\kappa)$ \\
\midrule
23   & 4.205 & 3.105 & 1.100 & 0.354 \\
53   & 3.707 & 3.105 & 0.601 & 0.194 \\
101  & 4.437 & 3.105 & 1.332 & 0.429 \\
199  & 4.592 & 3.105 & 1.486 & 0.479 \\
503  & 4.662 & 3.105 & 1.556 & 0.501 \\
1009 & 4.915 & 3.105 & 1.809 & 0.583 \\
\bottomrule
\end{tabular}
\end{center}
The maximum observed ratio is $\rho_{\max}=0.583<1$,
consistent with $(E_{\mathrm{rem}})$ for any $\rho\in(0.583,1)$.
Whether $\rho(\kappa)$ remains below $1$ for all $\kappa$ is
analytically OPEN.
\end{remark}

\begin{remark}[Route toward $(E_{\mathrm{rem}})$]
A stronger but more structured hypothesis $(E_{\mathrm{osc}})$ would
assert uniform cancellation of weighted prime-phasor partial sums:
\[
\left|\sum_{p\le\kappa} c_p^2\,
      \sin(\gamma_{k_1}\log p)\sin(\gamma_{k_2}\log p)
\right| \le C\cdot\frac{\Dburstvar}{\pi(\kappa)^A}
\]
for some $C>0$, $A\ge 1$, independent of $k_1,k_2$.
Such an estimate would still have to be combined
with a separate aggregation lemma controlling the
passage from these pairwise oscillatory sums to
$|\Dcross+\Dstream|$.
That implication is not proved here and is part
of the open problem; see Open Problem~\ref{op:weyl}.
The related hypothesis $(E_{\mathrm{pair}})$ controls pairwise
cancellation at a weaker level; its implication for $(E_{\mathrm{rem}})$
is also not established.
\end{remark}

\begin{remark}[Correction to earlier formulation]
The earlier claim ``$\Delta(\kappa)\to\Dburstvar$'' was incorrect:
$\Dburstvar \approx 3.11$ is a $\kappa$-invariant positive contribution,
not the limit. Numerically, $\Delta_\infty \approx 4.9 > \Dburstvar$.
It is observed to be a lower bound for $\Delta(\kappa)$ on the tested grid,
but this lower-bound property is not proved without additional assumptions.
\end{remark}

\begin{remark}[Numerical extrapolation]
On the tested grid $\kappa\in\{23,53,101,199,503,1009\}$, $\etaren(\kappa)$
remains positive and lies in $[0.669, 0.820]$. The data are compatible with a positive limiting value,
but no convergence theorem is proved.
\end{remark}

\begin{remark}[Connection to RH]
Under $(E_{\mathrm{rem}})$, Theorem~\ref{thm:T5} gives $\etaren>0$
at the stated reference parameters.
A possible future bridge toward RH would require an independent
transfer theorem connecting an appropriate limiting form of
$\etaren>0$ to zero exclusion off the critical line.
No such limiting theorem or transfer principle is established here.
\end{remark}

% -------------------------------------------------------
\section{Open Problems}\label{sec:open-problems}

\begin{openproblem}[Rank equality $\mathrm{rank}(\Tt)=\pi(\kappa)$]\label{op:rank-eq}
At benchmark parameters $(\kappa,\varepsilon,N)=(53,0.05,100)$, prove
analytically that $\{a_p\}_{p\leq\kappa}\subset\Hnull$ are linearly
independent --- equivalently that $\mathrm{rank}(\Tt)=\pi(\kappa)$
rather than the unconditional bound $\mathrm{rank}(\Tt)\leq\pi(\kappa)$
established in Proposition~\ref{prop:Tt-properties} via the at-most
inequality. Numerically confirmed: the spectral gap
$\mu_{\pi(\kappa)+1}/\mu_{\pi(\kappa)}\approx 3.7\cdot 10^{-13}$ at
benchmark parameters is consistent with linear independence and hence
with $\mathrm{rank}(\Tt)=\pi(\kappa)$ at machine precision.
\end{openproblem}

\begin{openproblem}[Proof of $(E_{\mathrm{rem}})$]\label{op:weyl}
Prove Assumption $(E_{\mathrm{rem}})$ (Definition~\ref{ass:Erem})
at fixed reference parameters $(\varepsilon,N)=(0.05,100)$:
that there exist $\rho<1$ and $\kappa_0$ such that
$|\Dcross(\kappa)+\Dstream(\kappa)|\le\rho\cdot\Dburstvar$
for all $\kappa\ge\kappa_0$.
Numerically $\rho_{\max}=0.583$ on the tested grid
$\kappa\in\{23,53,101,199,503,1009\}$; whether this remains below~$1$
for all $\kappa$ is open.

A structured approach via $(E_{\mathrm{osc}})$: prove
\[
\left|\sum_{p\le\kappa} c_p^2\,
      \sin(\gamma_{k_1}\log p)\sin(\gamma_{k_2}\log p)
\right| \le C\cdot\frac{\Dburstvar}{\pi(\kappa)^A}
\]
for some $C>0$, $A\ge 1$.
By itself, this estimate would \textbf{not} yet prove $(E_{\mathrm{rem}})$;
one would also need an aggregation lemma transferring the pairwise
oscillatory control to $|\Dcross+\Dstream|$.
The formulation and proof of such an aggregation lemma
are part of this open problem.
This approach is adjacent to classical questions on zero-ordinate
independence and prime phase distribution, but it is a separate
quantitative prime-phase cancellation problem.
No implication to or from the classical Linear Independence
conjecture is asserted here; for context see
\cite{Ingham42,RubinsteinSarnak94,FordMengZaharescu17}.
The related hypothesis $(E_{\mathrm{pair}})$ controls pairwise
cancellation at a weaker level; its implication for $(E_{\mathrm{rem}})$
is also not established.
\end{openproblem}

\begin{openproblem}[Analytic positivity of $\etaren$]\label{op:pos}
Prove $\etaren(\kappa)>0$ for all $\kappa\ge 2$ without assuming
$(E_{\mathrm{rem}})$ or any unverified hypothesis
(at fixed reference parameters $(\varepsilon,N)=(0.05,100)$).

Current best: Proposition~\ref{thm:EXPLICIT} gives a PNT-based bound
for raw prime exponential sums $\sum_{p\le x}\sin(\gamma_k\log p)$
at fixed $\gamma_k$. It does not yet bound the spectral projections
$|\langle u_j,\tilde{c}\rangle|$, where $\tilde{c}=\Da^{1/2}c/\sqrt{\Estr}$
is the normalised vector in the Rayleigh quotient identity of \S\,5.1.
A missing transfer lemma would be needed to pass from the prime-phasor
bounds to the normalised spectral energy control
\[
\sum_j\lambda_j|\langle u_j,\tilde{c}\rangle|^2 < 1.
\]
Equivalently, in unnormalised form:
\[
\sum_j\lambda_j|\langle u_j,\Da^{1/2}c\rangle|^2 < \Estr.
\]
\end{openproblem}

\begin{openproblem}[Circularity-free spectral function]\label{op:hp}
Two separate questions.

Here ``circularity-free'' means that the functional calculus rule
is fixed independently of any OLS fit, postulated HP spectrum, or
post hoc matching to the target ordinates; the ordinates entering
$\Tt$ through $\Phi$ remain external inputs.

\emph{(a) Spectral function of the fixed $\Tt$:}
For fixed $N$, does any circularity-free function $f(\Tt_{\kappa,N})$
have its first $N$ spectral values converge, after a specified matching,
to $\{\gamma_1,\ldots,\gamma_N\}$ as $\kappa\to\infty$?
The natural candidate $\Wone = \CT\cdot\Tt^+$ fails: $r_2$ falls to $0.16$ on the tested grid
(\S\,3.4). The observed obstruction is consistent with Gaussian damping;
no structural obstruction theorem is proved here.

\emph{(b) Model-variant question (separate from (a)):}
If the Gaussian damping $e^{-\varepsilon^2\gamma_k^2/2}$ in $\Phi$ were
replaced by a slower decay $e^{-\varepsilon\gamma_k}$, could the
resulting modified operator $\Tt'$ support a circularity-free
spectral function? This would define a different operator family
and is not a spectral function of the present $\Tt$.
\end{openproblem}

\begin{openproblem}[Arithmetic constant $\Ceta$]\label{op:ceta}
Determine the analytic formula for the observed ratio
$R_\eta(\kappa)=\lambda_{\max,\mathrm{ren}}/\pi(\kappa)\approx 0.39$
(denoted $\Ceta$; see Definition~\ref{def:constants}).
The candidate $\Ceta = 1/I_N(\varepsilon)$ is falsified.
The correct representation involves a harmonic mean
(analytically open).
\end{openproblem}

\begin{openproblem}[Asymptotics of $\CT$]\label{op:ct}
Does $\CT(\kappa)/\pi(\kappa)$ converge as $\kappa\to\infty$?
Numerically: values $2.49, 2.21, 2.17, 1.94, 1.89$ (decreasing) ---
the limit, if it exists, is analytically undetermined.
\end{openproblem}

\begin{openproblem}[Monotone bound from prime exponential-sum control]\label{op:mono}
Does $\etaren > 0$ follow from PNT-based prime exponential-sum
control, with the zero ordinates $\gamma_k$ treated as fixed
external parameters (at fixed $(\varepsilon,N)=(0.05,100)$)?

Evidence: $\Dburstvar\approx 3.11 > 0$ depends only on primes
$p\leq 7$ and is $\kappa$-invariant.
Analytically OPEN.
\end{openproblem}

\begin{remark}[Connection to RH]
Under $(E_{\mathrm{rem}})$, Theorem~\ref{thm:T5} gives $\etaren > 0$.
A possible future bridge toward RH would require an independent
transfer theorem connecting an appropriate limiting form of
$\etaren>0$ to zero exclusion off the critical line.
No such limiting theorem or transfer principle is established here.
\end{remark}

% -------------------------------------------------------
\section*{Non-Circularity}
\label{sec:nocirc}

We document explicitly that the results of this paper do not rely on
the objects they are directed toward. The following holds unconditionally
throughout \S\S\,2--5.

\textbf{No RH as input.}
At no point is the Riemann Hypothesis assumed. The zero ordinates
$\gamma_k$ enter the construction of $\Phi$ as numerically specified
inputs, and no property of their distribution requiring the Riemann
Hypothesis is assumed or used. In Theorem~\ref{thm:T5}, Assumption $(E_{\mathrm{rem}})$ is
explicitly labelled CONDITIONAL; it concerns the energy decomposition
and is not a consequence of RH.

\textbf{No GUE, Montgomery, or Hilbert--P\'olya hypothesis.}
The Tehrani operator $\Tt = \Phi\circ\Phi^*$ is a deterministic,
explicitly computable finite-dimensional operator, constructed from
the prime resonance vectors $\{a_p\}_{p\le\kappa}$. It is not a random
matrix, not a statistical ensemble, and not a Hilbert--P\'olya
operator. No GUE conjecture, no Montgomery pair correlation conjecture,
and no random-matrix-theoretic assumption is used or implied anywhere
in \S\S\,2--5. The eigenvector-localization pattern reported in
\S\,3 is a numerical observation, not a postulate.

\textbf{Status separation maintained.}
The zero ordinates $\gamma_k$ appear as inputs to the construction
of $\Phi$, never as conclusions. The eigenvalues $\mu_j$ of $\Tt$
are derived from $\Phi$ and are not identified with the $\gamma_k$.

\textbf{Note on $\Wone$.}
The inverse operator $\Wone = \CT\cdot\Tt^+$ does not assume RH or
the Hilbert--P\'olya postulate; however, its normalization constant
$\CT$ is calibrated against the zero ordinates via OLS (\S\,3.3).
$\Wone$ is therefore not fully input-free. $\Tt$ itself is
deterministic and explicitly computable.

% -------------------------------------------------------
\section*{Acknowledgments}

The author thanks the reviewers and correspondents whose
comments improved the presentation of this paper.

% -------------------------------------------------------
\section*{Call for Feedback}

This paper is part of an open research initiative on prime--zero
coupling and explicit-formula operators.
All papers are freely available on Zenodo under CC BY 4.0, and the
accompanying code is hosted on GitHub.

{\raggedright
Zenodo:\\
\url{https://zenodo.org/search?q=metadata.creators.person_or_org.name\%3A\%22Tehrani\%2C\%20Ulrich\%22}\\
GitHub:\\
\url{https://github.com/utehrani/analysislab-nt}
\par}

The author welcomes critical feedback, corrections, and collaboration.
In particular: analytical approaches to Assumption $(E_{\mathrm{rem}})$
(remainder control for $|\Dcross+\Dstream|$,
Open Problem~\ref{op:weyl}) and its relation to the
Linear Independence conjecture and prime exponential sums,
connections to existing results on exponential sums
over primes and zeros of $\zeta(s)$,
and any errors or gaps in the present arguments.

% -------------------------------------------------------
% Bibliography — normative format for the series
% Series papers: Concept-DOI only, no version numbers.
% External refs: standard bibliographic format.
% -------------------------------------------------------
\newpage
\begin{thebibliography}{99}

\bibitem{BerryKeating}
M.V.~Berry and J.P.~Keating,
``$H=xp$ and the Riemann zeros'',
in: J.P.~Keating, D.E.~Khmelnitskii, I.V.~Lerner (eds.),
\textit{Supersymmetry and Trace Formulae: Chaos and Disorder},
NATO ASI Series B: Physics, vol.~370,
Plenum Press, New York, 1999, pp.~355--367.

\bibitem{Bombieri}
E.~Bombieri,
``Problems of the Millennium: The Riemann Hypothesis'',
Clay Mathematics Institute, 2000.

\bibitem{BostConnes}
J.-B.~Bost and A.~Connes,
``Hecke algebras, type III factors and phase transitions
with spontaneous symmetry breaking in number theory'',
\emph{Selecta Mathematica, New Series}\ \textbf{1} (1995), no.~3, 411--457.

\bibitem{BurnolForum04}
J.-F.~Burnol,
``Two complete and minimal systems associated with
the zeros of the Riemann zeta function'',
\emph{J. Th\'eor. Nombres Bordeaux}\ \textbf{16} (2004), no.~1, 65--94.

\bibitem{CCM25}
A.~Connes, C.~Consani, and H.~Moscovici,
``Zeta spectral triples'',
arXiv:2511.22755 (2025).

\bibitem{Connes}
A.~Connes,
``Trace formula in noncommutative geometry and the zeros
of the Riemann zeta function'',
\emph{Selecta Math.}\ \textbf{5} (1999), no.~1, 29--106.

\bibitem{ConnesConsani21}
A.~Connes and C.~Consani,
``Weil positivity and trace formula, the archimedean place'',
\emph{Sel. Math. (N.S.)}\ \textbf{27} (2021), Paper No.~77.

\bibitem{EndresSteiner10}
S.~Endres and F.~Steiner,
``The Berry--Keating operator on $L^2(\mathbb{R}_{>},\mathrm{d}x)$
and on compact quantum graphs with general self-adjoint realizations'',
\emph{J. Phys. A: Math. Theor.}\ \textbf{43} (2010), 095204.

\bibitem{FordMengZaharescu17}
K.~Ford, X.~Meng, and A.~Zaharescu,
``Zeros of the Riemann zeta-function on the critical line
and the distribution of the argument'',
\emph{Bull. London Math. Soc.}\ \textbf{49} (2017), 1--9.

\bibitem{Ingham42}
A.E.~Ingham,
``On two conjectures in the theory of numbers'',
\emph{Amer. J. Math.}\ \textbf{64} (1942), 313--319.

\bibitem{IwaniecKowalski}
H.~Iwaniec and E.~Kowalski,
\emph{Analytic Number Theory},
AMS Colloquium Publications~53, 2004.

\bibitem{RubinsteinSarnak94}
M.~Rubinstein and P.~Sarnak,
``Chebyshev's bias'',
\emph{Experimental Math.}\ \textbf{3} (1994), 173--197.

\bibitem{SchumayerHutchinson11}
D.~Schumayer and D.A.W.~Hutchinson,
``Colloquium: Physics of the Riemann hypothesis'',
\emph{Rev. Mod. Phys.}\ \textbf{83} (2011), 307--330.

\bibitem{SierraRodriguez11}
G.~Sierra and J.~Rodr\'iguez-Laguna,
``$H=xp$ model revisited and the Riemann zeros'',
\emph{Phys. Rev. Lett.}\ \textbf{106} (2011), 200201.

\bibitem{SierraTownsend}
G.~Sierra and P.K.~Townsend,
``Landau levels and Riemann zeros'',
\emph{Physical Review Letters}\ \textbf{101} (2008), 110201.

\bibitem{Turan}
P.~Tur\'an,
``On some approximative Dirichlet-polynomials in the theory of
the zeta-function of Riemann'',
\emph{Danske Videnskabernes Selskab, Mat.-fys.\ Medd.}\
\textbf{24} (1948), no.~17, 3--36 (36~pp.).

% Series papers — Concept-DOI only, no version numbers
% Prior-series citations only; no forward-citations
\bibitem{Paper1}
U.~Tehrani,
``A Curvature Decomposition of the Explicit Formula
for the Riemann Zeta Function'',
Zenodo, \url{https://doi.org/10.5281/zenodo.19025598}, 2026.

\bibitem{Paper2}
U.~Tehrani,
``From Local Curvature to the Weil Functional: An Explicit Construction'',
Zenodo, \url{https://doi.org/10.5281/zenodo.19106992}, 2026.

\bibitem{Paper3}
U.~Tehrani,
``A Finite-Cutoff Hilbert-Space Model
for Prime--Zero Energy Structure'',
Zenodo, \url{https://doi.org/10.5281/zenodo.19307989}, 2026.

\end{thebibliography}

% --------------------------------------------------------
% Footer — normative format
% --------------------------------------------------------
\enlargethispage{3\baselineskip}
\vspace*{\fill}
\begin{minipage}{\linewidth}
\rule{\linewidth}{0.4pt}\smallskip\\
\small\emph{Paper~4 v31 \textperiodcentered\ August~2026
\textperiodcentered\ Pauca sed matura}
\end{minipage}

\end{document}
